\section{The objects being counted}\label{sec:definitions}
For deletion arguments, \(n\geq1\). The extremal functions themselves
are also defined at \(n=0\); that initial value is harmless in limits.
\begin{definition}[Orders, profiles and domains]\label{def:semantic}
An order on a finite set \(X\) is a strict total order, written as a list
from best to worst. A domain is a set \(D\) of orders on the same \(X\).
For a profile \((\sigma_1,\ldots,\sigma_m)\), a strict majority puts
\(a\) above \(b\) when more than \(m/2\) voters do so. A Condorcet
domain is a domain for which this majority relation is transitive for
every odd positive \(m\) and every profile from \(D\). Repeated orders
in a profile are allowed. Write \(f(n)\) for the maximum possible
\(|D|\) on \(n\) alternatives.
\end{definition}
\lean{\path{Condorcet.IsStrictTotal}, \path{Condorcet.Order},
\path{Condorcet.Domain}, \path{Condorcet.majority},
\path{Condorcet.IsCondorcet}, \path{Condorcet.f}.
An order is a Boolean relation with irreflexivity, transitivity and
totality; asymmetry follows. A domain is a \texttt{Finset}, so its
members are distinct. Profiles are functions \(\mathrm{Fin}(m)\to
\mathrm{Order}(X)\), with no injectivity requirement. The definition
of \(f\) takes a finite supremum over exactly these semantic domains.
The supplied correspondence source restates these three semantic interfaces in
\path{TutorialCorrespondence.condorcet_definition},
\path{peak_pit_definition} and \path{triple_characterization}.}

\subsection{Triple laws connect voting to combinatorics}
On a triple \(\{x,y,z\}\), the law \(xN_i\) says that \(x\) never
occupies position \(i\), with \(i=1,2,3\) meaning first, middle, last.
An extreme law is \(xN_1\) or \(xN_3\).
\begin{definition}[Peak-pit]\label{def:peak}
A peak-pit domain has a common extreme law on each triple, with the
choice of alternative and forbidden position allowed to vary between
triples. Let \(h(n)\) be its maximum size. No common axis is required.
\end{definition}
\lean{\path{Condorcet.Never} uses positions \(0,1,2\), respectively.
\path{Condorcet.ValueRestricted} means a common never-law per triple;
\path{ExtremeRestricted} requires the position to differ from \(1\)
in this zero-based convention. \path{IsPeakPit} is explicitly
\path{IsCondorcet} together with \path{ExtremeRestricted};
\path{Condorcet.h} takes the corresponding maximum.}

\begin{theorem}[Triple characterisation]\label{thm:triple}
A domain is Condorcet if and only if each triple has a common never-law.
In particular, a common extreme law per triple suffices.
\end{theorem}
\begin{proof}
On a fixed triple, separate the six orders into
\[
 \{abc,bca,cab\},\qquad\{acb,cba,bac\}.
\]
The three orders in either set, used once each as a three-voter profile,
produce a majority cycle. A Condorcet domain must therefore omit at
least one order from each set. There are nine choices of the two
omitted orders. In every choice, some alternative occupies the same
position in both omitted orders; the remaining four orders satisfy
that never-law. To see the nine cases without a search, fix the omitted
order in the first set to be \(abc\) by cyclic relabelling. Pairing
it with \(acb,cba,bac\) gives respectively \(aN_1,bN_2,cN_3\).
Undo the relabelling for the other two first-set choices.

Conversely, suppose all voters satisfy a fixed never-law on a triple.
For \(xN_2\), each voter compares \(x\) with \(y,z\) in the same
direction, so the two majority comparisons agree and cannot form a
cycle. For \(xN_1\), a voter with \(x>y\) must have \(z>x>y\).
Thus a majority for \(x>y\) forces a majority for \(z>y\), ruling
out the cycle \(x>y>z>x\). For the reverse cycle, interchange \(y,z\).
The case \(xN_3\) is the reversal of this argument.
An odd profile has no majority ties. Its majority relation is a
tournament, and a non-transitive tournament has a directed triangle:
if \(a>b\), \(b>c\), but not \(a>c\), then \(c>a\). Every such
triangle is excluded by its common law, proving transitivity.
\end{proof}
\lean{\path{Condorcet.isCondorcet_iff_valueRestricted} is the exact
semantic equivalence. Its two directions are
\path{IsCondorcet.valueRestricted} and
\path{ValueRestricted.isCondorcet}. The finite triple step is
\path{Condorcet.Pattern.finite_characterization}; its finite
formulation uses the six Boolean comparison patterns. The proof above
is a human proof of the same equivalence, not a literal transcript of
Lean's finite case split.}

\subsection{Maximum, maximal and full are three different notions}
\begin{definition}[Fullness and the simpler class]\label{def:full}
A domain is maximal if no further order can be added while retaining
the Condorcet property. It has maximum size if its size equals \(f(n)\).
Choose one never-law on each triple. Its full solution set is called a
full domain. Define \(\F_n\) to consist of nonempty full solution sets
of one extreme law per triple together with any set of fixed pair
comparisons. A fixed comparison may be absent or may require either
orientation. It does not impose a common axis.
\end{definition}
\lean{\path{Condorcet.IsMaximal} includes both the semantic Condorcet
property and inclusion-maximality. \path{Selection},
\path{HoldsLaw}, \path{fullDomain} and \path{IsFull} formalise full
general domains. \path{ExtremeSelection}, \path{FixedPairs},
\path{extremeFullDomain} and \path{IsExtremeFull} describe the full
extreme-law-plus-pairs solution set, which may be empty.
\path{Condorcet.InF} adds nonemptiness and is exactly our \(\F_n\).
These two predicates must not be interchanged.}

\begin{lemma}[Completion]\label{lem:completion}
A maximum general domain is maximal, and a maximal general domain is
full. A maximum peak-pit domain is a full solution set of extreme
laws, and hence belongs to \(\F_n\).
\end{lemma}
\begin{proof}
Adding an order to a maximum domain would increase its size, proving
maximality. For a maximal domain choose a common law on each triple
using Theorem~\ref{thm:triple}, and include every order satisfying the
chosen laws. The enlargement is still Condorcet, so maximality makes
it equal to the original domain. For a maximum peak-pit domain, make
the choices extreme. Its completion remains peak-pit, so maximum size
again forces equality. It is nonempty, since a singleton is peak-pit.
There is no need for it to be maximal among \emph{all} Condorcet domains.
\end{proof}
\lean{\path{Condorcet.maximum_isMaximal}, \path{IsMaximal.isFull},
\path{ExtremeRestricted.completion},
\path{maximum_peakPit_isExtremeFull},
\path{maximum_peakPit_fullClass}. The last gives
\path{IsExtremeFull}; nonemptiness is supplied where \path{InF}
is used. The extremal functions are actual attained maxima by
\path{exists_maximum_general} and \path{exists_maximum_peakPit}.}
